%!TEX program = lualatex
\documentclass[ngerman,headheight=70pt]{scrartcl}
\usepackage[ngerman]{babel}
\usepackage{blindtext}
\usepackage[utf8]{luainputenc}
\usepackage{scrlayer-scrpage}
\usepackage[T1]{fontenc}
\usepackage{fontspec}
\usepackage{multicol}
\usepackage{enumitem}
\usepackage{epigraph}
\usepackage[german=quotes]{csquotes}
\usepackage[bottom=2cm,footskip=8mm]{geometry}
\setmainfont{Source Sans Pro}
\MakeOuterQuote{"}

\setlength{\parindent}{0pt}
\setlength{\parskip}{5pt}

\pagestyle{plain.scrheadings}
\thispagestyle{scrheadings}
\lohead{\vspace*{-1.5cm}\includegraphics[scale=0.4]{../images/stupa_logo_rot.pdf}}
\rohead{\\Präsidium des\\Studierendenparlaments}

\lofoot[Universität Hamburg · Präsidium des Studierendenparlaments\\\hrule Von-Melle-Park 5 · D-20146 Hamburg]{Universität Hamburg · Präsidium des Studierendenparlaments\\\hrule Von-Melle-Park 5 · D-20146 Hamburg}
\cofoot[]{}
\cohead[\thepage]{}

\renewcommand{\thesubsection}{TOP \arabic{subsection}}
\renewcommand{\thesubsubsection}{\arabic{subsubsection}.}
\setcounter{subsection}{-1}

\begin{document}
    UHH · StuPa-Präsidium · Von-Melle-Park 5 · D-20146 Hamburg

    \section*{Beschluss des Studierendenparlamentes vom 28. April 2016}
    \subsection*{Studienvorbereitung für Geflüchtete ausbauen!}

    Mit dem Programm \#uhh hilft ist es im gruppenübergreifenden Engagement vieler
    Unimitglieder gelungen, ein Vorbereitungsstudium für Flüchtlinge auf den Weg
    zu bringen.

    Das Studierendenparlament begrüßt das Engagement aller in diesem Bereich Aktiven.

    Mit dem Bemühen, das Programm für alle geflüchteten Studieninteressierten,
    die eine Hochschulzugangsberechtigung haben, zu öffnen, setzt die Uni ein
    Signal: gegen die Auslese auch von Flüchtlingen nach ihrer "Verwertbarkeit".
    Vielmehr geht es um zügige gesellschaftliche Partizipation auf hohem Niveau
    für möglichst viele Menschen. Und es geht darum, den Elendsbedingungen in den
    Flüchtlingsunterkünften entgegenzuwirken. Deshalb muß das Programm ausgebaut
    werden!

    Eklatant ist, daß in Folge von Kürzungsorgien und Fahrlässigkeit die Uni
    nicht in der Lage ist, bedarfsgerecht Kurse in "Deutsch als Fremdsprache"
    anzubieten. Es fehlen Intensivkurse auf A1/A2 und B1 Niveau, für Geflüchtete
    sowie für Austausch- und Programmstudierende. Zur Zeit erhalten im derzeitigen
    Vorbereitungsprogramm nur Teilnehmer*innen, die schon ein B1-Sprachzertifikat
    haben, Deutschkurse. Die Gruppen geflüchteter Studieninteressierter, die keinen
    Anspruch auf staatlich finanzierte Deutschkurse auf Anfängerniveau haben und
    -- nach Willen der Bundes- und Landesregierung -- jahrelang ohne sinnvolle
    Beschäftigung auf ihr Asylverfahren warten sollen, sind damit der Möglichkeit,
    sich auf ein Studium adäquat vorzubereiten, beraubt. Das trifft besonders
    Menschen aus sogenannten "sicheren Herkunftsstaaten" wie Afghanistan und aus
    dem Maghreb. Politisch wird damit zwischen den Flüchtlingen (mit oder ohne
    sogenannte "sichere Bleibeperspektive") zu spalten versucht.


    Das Studierendenparlament lehnt diese Politik der Auslese, Abwertung,
    Abschreckung und Abschiebung ab.

    Die Universität hat mit \#uhh hilft die Möglichkeit, dieser Ungerechtigkeit
    beispielgebend entgegenzuwirken. Das Studierendenparlament fordert daher die
    Universitätsleitung dazu auf:

    \begin{itemize}
        \item zügig das Angebot an Deutschintensivkursen für Anfäger*innen
        bedarfsgerecht für Programm-Teilnehmer*innen und Universitätsmitglieder
        auszubauen.
        \item Den DaF-Dozent*innen angemessene, längerfristige Arbeitsverträge
        zu ermöglichen.
        \item Dadurch einerseits mehr Dozent*innen zu gewinnen, aber ihnen auch
        bessere Arbeitsbedingungen zu schaffen, die ein gutes Lehr-Lern-Verhältnis
        (auch außerhalb der Kurse), persönliche Beratung und die gemeinsame
        Reflektion der Arbeit (besonders mit Geflüchteten) ermöglichen.
        \item Zur Organisation der Sprachkurse dem Büro der Flüchtlingsbeauftragten
        die erforderlichen Mittel für Personal zur Verfügung zu stellen.
    \end{itemize}

    Diese Maßnahmen sind akut notwendig und können zudem dem erforderlichen
    Wiederaufbau eines wissenschaftlichen Arbeitsbereichs Deutsch als Fremdsprache
    den Weg bereiten.
    Darüber hinaus fordert das Stupa die Bürgerschaft auf, den Hamburger
    Hochschulen für den bedarfsgerechten Ausbau von DaF und die Durchführung der
    Studien(vorbereitungs)programme für Geflüchtete dauerhaft bedarfsgerechte Mittel
    zusätzlich zuzuweisen.

    %\vspace{0.4cm}
    \textbf{Das StuPa-Präsidium wird gebeten, diesen Beschluss dem Unipräsidium
    zuzuleiten.}
\end{document}
